\documentclass{article}
\usepackage[legalpaper, portrait, margin=0.5in]{geometry}
\usepackage{amsmath, amssymb, amsthm}
\usepackage{hyperref}
\usepackage{enumitem}
\usepackage{microtype}
\usepackage{array}
\usepackage{booktabs}
\hypersetup{ colorlinks=true, linkcolor=blue, urlcolor=blue }

\newcommand{\at}[1]{\text{#1}}
\newcommand{\cl}[1]{\{#1\}^+}

\title{CSC 370 PaperTrade, Lossless Join Proof}
\author{Dryden Bryson}
\date{\today}

\begin{document}

\maketitle
\newpage

\section*{Overview}

Each of our four relations (\texttt{user}, \texttt{stock}, \texttt{transaction} and \texttt{holds}) is already in BCNF (see the \emph{Functional Dependencies and Candidate Keys} document). Losslessness is a property of a \emph{decomposition}, not of a single relation. Our four relations are the result of decomposing the single table we would have without normalization: the \textbf{universal relation} $U$, which holds every attribute of the database in one row per trade.

This document:
\begin{enumerate}[leftmargin=*]
  \item defines $U$ and its FDs, finds its candidate key, and shows that it violates BCNF, using anomalies from \texttt{sql/seed.sql};
  \item decomposes $U$ with the BCNF decomposition algorithm, which ends in exactly our four relations, and proves each step lossless;
  \item confirms the whole four-way decomposition is lossless with the chase test;
  \item gives a lossy decomposition of \texttt{holds} as a counterexample, with spurious tuples taken from the seed data.
\end{enumerate}

\section{The Universal Relation}

$U$ stores one row per trade. Each row also carries the trader's account details, the stock's details, and the trader's current position in that stock. \texttt{transaction} and \texttt{holds} both have a \texttt{quantity} attribute, so we rename them $\at{t\_quantity}$ (shares in this trade) and $\at{h\_quantity}$ (shares currently held).
\[
  U(\at{transaction\_id}, \at{user\_id}, \at{name}, \at{email}, \at{ticker}, \at{company\_name}, \at{opening\_price}, \at{closing\_price}, \at{exchange},
\]
\[
  \at{transaction\_type}, \at{t\_quantity}, \at{price}, \at{time}, \at{h\_quantity}, \at{average\_price})
\]

To keep the expressions readable, we group attributes that always appear together:
\begin{align*}
  N &= \{\at{name}, \at{email}\} && \text{(account details)} \\
  S &= \{\at{company\_name}, \at{opening\_price}, \at{closing\_price}, \at{exchange}\} && \text{(stock details)} \\
  T &= \{\at{transaction\_type}, \at{t\_quantity}, \at{price}, \at{time}\} && \text{(trade details)} \\
  P &= \{\at{h\_quantity}, \at{average\_price}\} && \text{(position details)}
\end{align*}
\[
  U = \{\at{transaction\_id}, \at{user\_id}, \at{ticker}\} \cup N \cup S \cup T \cup P
\]

\subsection*{Functional Dependencies $F$}

These are the FDs of the four relations, from the \emph{Functional Dependencies and Candidate Keys} document:
\begin{itemize}[leftmargin=*]
  \item $FD_1$: $\at{user\_id} \rightarrow N$
  \item $FD_2$: $\at{email} \rightarrow \at{user\_id}$
  \item $FD_3$: $\at{ticker} \rightarrow S$
  \item $FD_4$: $\at{transaction\_id} \rightarrow \at{user\_id}, \at{ticker}, T$
  \item $FD_5$: $\at{user\_id}, \at{ticker} \rightarrow P$
\end{itemize}

\subsection*{Candidate Key of $U$}

\begin{align*}
  \cl{\at{transaction\_id}}
    &\supseteq \{\at{transaction\_id}\} && \text{(reflexivity)} \\
    &\supseteq \{\at{transaction\_id}, \at{user\_id}, \at{ticker}\} \cup T && \text{(applying } FD_4\text{)} \\
    &\supseteq \{\at{transaction\_id}, \at{user\_id}, \at{ticker}\} \cup T \cup N && \text{(applying } FD_1\text{, transitivity)} \\
    &\supseteq \{\at{transaction\_id}, \at{user\_id}, \at{ticker}\} \cup T \cup N \cup S && \text{(applying } FD_3\text{, transitivity)} \\
    &\supseteq \{\at{transaction\_id}, \at{user\_id}, \at{ticker}\} \cup T \cup N \cup S \cup P && \text{(applying } FD_5\text{, transitivity)} \\
    &= U
\end{align*}

$\{\at{transaction\_id}\}$ is a superkey, and its only proper subset $\emptyset$ has $\emptyset^+ = \emptyset$, so it is a candidate key. No FD has $\at{transaction\_id}$ on its right-hand side. The only way to derive it is reflexivity, so every superkey must contain $\at{transaction\_id}$. $\{\at{transaction\_id}\}$ is therefore the \textbf{only} candidate key of $U$.

\subsection*{BCNF Violations}

$U$ is in BCNF only if the left-hand side of every non-trivial FD is a superkey. Four of the five FDs fail this:
\begin{center}
\begin{tabular}{l l l}
  \toprule
  \textbf{FD} & \textbf{Closure of left-hand side in $U$} & \textbf{Superkey?} \\
  \midrule
  $FD_1$ & $\cl{\at{user\_id}} = \{\at{user\_id}\} \cup N$ & No \\
  $FD_2$ & $\cl{\at{email}} = \{\at{user\_id}\} \cup N$ & No \\
  $FD_3$ & $\cl{\at{ticker}} = \{\at{ticker}\} \cup S$ & No \\
  $FD_4$ & $\cl{\at{transaction\_id}} = U$ & Yes \\
  $FD_5$ & $\cl{\at{user\_id}, \at{ticker}} = \{\at{user\_id}, \at{ticker}\} \cup N \cup S \cup P$ & No \\
  \bottomrule
\end{tabular}
\end{center}

\subsection*{Anomalies in $U$}

Each violation causes anomalies, which the seed data shows directly:
\begin{itemize}[leftmargin=*]
  \item \textbf{Insertion:} every row of $U$ needs a $\at{transaction\_id}$. So $U$ cannot store user 9 (Isla MacDonald, who has no trades yet), or the stocks CNR and ATD, which have never been traded.
  \item \textbf{Update:} AAPL appears in 9 trades, so its closing price is stored 9 times. Changing it means updating all 9 rows, and missing one leaves $U$ inconsistent ($FD_3$).
  \item \textbf{Deletion:} user 7 (Grace Liu) has 4 trades. Deleting those 4 rows also deletes Grace's name and email ($FD_1$).
\end{itemize}

\section{BCNF Decomposition of $U$}

\textbf{Lossless join test.} A decomposition of $R$ into $R_1$ and $R_2$ is lossless if and only if
\[
  (R_1 \cap R_2) \rightarrow R_1 \quad \text{or} \quad (R_1 \cap R_2) \rightarrow R_2 \quad \text{is in } F^+
\]
That is, the shared attributes form a superkey of at least one side.

Each step of the algorithm picks a violation $X \rightarrow Y$ and splits $R$ into $R_1 = X^+$ and $R_2 = (R \setminus X^+) \cup X$. Then $R_1 \cap R_2 = X$, and $X \rightarrow X^+ = R_1$ by the definition of closure, so each step passes the test.

\subsection*{Step 1: Remove $FD_1$}

\begin{itemize}[leftmargin=*]
  \item \textbf{Violation:} $FD_1$: $\at{user\_id} \rightarrow N$, with $\cl{\at{user\_id}} = \{\at{user\_id}\} \cup N$.
  \item \textbf{Split:}
  \begin{align*}
    R_1 &= \{\at{user\_id}\} \cup N = \{\at{user\_id}, \at{name}, \at{email}\} \implies \textbf{\texttt{user}} \\
    U_1 &= (U \setminus N) \cup \{\at{user\_id}\} = \{\at{transaction\_id}, \at{user\_id}, \at{ticker}\} \cup S \cup T \cup P
  \end{align*}
  \item \textbf{Lossless:} $R_1 \cap U_1 = \{\at{user\_id}\}$, and $\cl{\at{user\_id}} \supseteq R_1$. So $\{\at{user\_id}\}$ is a superkey of $R_1$ and the split is lossless.
  \item $FD_2$ ($\at{email} \rightarrow \at{user\_id}$) now lies entirely inside \texttt{user}, where $\at{email}$ is a candidate key. So it is no longer a violation.
\end{itemize}

\subsection*{Step 2: Remove $FD_3$}

The FDs projected onto $U_1$ include $FD_3$, $FD_4$ and $FD_5$, since all their attributes lie in $U_1$. $\{\at{transaction\_id}\}$ is still the only key of $U_1$, so $FD_3$ still violates BCNF.
\begin{itemize}[leftmargin=*]
  \item \textbf{Violation:} $FD_3$: $\at{ticker} \rightarrow S$, with $\cl{\at{ticker}} = \{\at{ticker}\} \cup S$.
  \item \textbf{Split:}
  \begin{align*}
    R_2 &= \{\at{ticker}\} \cup S \implies \textbf{\texttt{stock}} \\
    U_2 &= (U_1 \setminus S) \cup \{\at{ticker}\} = \{\at{transaction\_id}, \at{user\_id}, \at{ticker}\} \cup T \cup P
  \end{align*}
  \item \textbf{Lossless:} $R_2 \cap U_2 = \{\at{ticker}\}$, and $\cl{\at{ticker}} \supseteq R_2$. So $\{\at{ticker}\}$ is a superkey of $R_2$ and the split is lossless.
\end{itemize}

\subsection*{Step 3: Remove $FD_5$}

The FDs projected onto $U_2$ include $FD_4$ and $FD_5$. In $U_2$, $\cl{\at{user\_id}, \at{ticker}} \cap U_2 = \{\at{user\_id}, \at{ticker}\} \cup P \neq U_2$, because it does not contain $\at{transaction\_id}$. So $FD_5$ still violates BCNF.
\begin{itemize}[leftmargin=*]
  \item \textbf{Violation:} $FD_5$: $\at{user\_id}, \at{ticker} \rightarrow P$.
  \item \textbf{Split:}
  \begin{align*}
    R_3 &= \{\at{user\_id}, \at{ticker}\} \cup P = \{\at{user\_id}, \at{ticker}, \at{h\_quantity}, \at{average\_price}\} \implies \textbf{\texttt{holds}} \\
    R_4 &= (U_2 \setminus P) = \{\at{transaction\_id}, \at{user\_id}, \at{ticker}\} \cup T \implies \textbf{\texttt{transaction}}
  \end{align*}
  \item \textbf{Lossless:} $R_3 \cap R_4 = \{\at{user\_id}, \at{ticker}\}$, and $\cl{\at{user\_id}, \at{ticker}} \supseteq R_3$. So $\{\at{user\_id}, \at{ticker}\}$ is a superkey of $R_3$ and the split is lossless.
\end{itemize}

\subsection*{Result}

The algorithm stops, because each of the four resulting relations is in BCNF:
\begin{center}
\begin{tabular}{l l l}
  \toprule
  \textbf{Relation} & \textbf{Attributes} & \textbf{Projected FDs} \\
  \midrule
  \texttt{user} & $\at{user\_id}, \at{name}, \at{email}$ & $FD_1$, $FD_2$ \\
  \texttt{stock} & $\at{ticker}, \at{company\_name}, \at{opening\_price}, \at{closing\_price}, \at{exchange}$ & $FD_3$ \\
  \texttt{holds} & $\at{user\_id}, \at{ticker}, \at{h\_quantity}, \at{average\_price}$ & $FD_5$ \\
  \texttt{transaction} & $\at{transaction\_id}, \at{user\_id}, \at{ticker}, \at{transaction\_type}, \at{t\_quantity}, \at{price}, \at{time}$ & $FD_4$ \\
  \bottomrule
\end{tabular}
\end{center}
This is exactly the schema in \texttt{sql/ddl.sql}. Every FD of $F$ lies inside one relation, so the decomposition also preserves dependencies.

Each binary step is lossless, so the whole decomposition is lossless. If $R = R_1 \bowtie U_1$ and $U_1 = R_2 \bowtie U_2$, then $R = R_1 \bowtie R_2 \bowtie U_2$, and so on for every step:
\[
  \pi_{\texttt{user}}(U) \bowtie \pi_{\texttt{stock}}(U) \bowtie \pi_{\texttt{holds}}(U) \bowtie \pi_{\texttt{transaction}}(U) = U
\]

\section{Chase Test of the Four-Way Decomposition}

The chase checks the whole decomposition at once. Build a tableau with one row per relation. Put the distinguished symbol $a$ in each column the relation contains, and a subscripted $b_{ij}$ (row $i$, column $j$) everywhere else. Each group $N$, $S$, $T$ and $P$ can be one column, because each group lies wholly inside or wholly outside every relation, and every FD's right-hand side contains whole groups.

\textbf{Initial tableau:}
\[
\begin{array}{l|ccccccc}
  & \at{transaction\_id} & \at{user\_id} & N & \at{ticker} & S & T & P \\
  \hline
  \texttt{user}        & b_{11} & a      & a      & b_{14} & b_{15} & b_{16} & b_{17} \\
  \texttt{stock}       & b_{21} & b_{22} & b_{23} & a      & a      & b_{26} & b_{27} \\
  \texttt{holds}       & b_{31} & a      & b_{33} & a      & b_{35} & b_{36} & a      \\
  \texttt{transaction} & a      & a      & b_{43} & a      & b_{45} & a      & b_{47} \\
\end{array}
\]

\textbf{Applying the FDs:}
\begin{enumerate}[leftmargin=*]
  \item $FD_1$ ($\at{user\_id} \rightarrow N$): rows \texttt{user}, \texttt{holds} and \texttt{transaction} agree on $\at{user\_id} = a$, and \texttt{user} has $N = a$. Set $b_{33} := a$ and $b_{43} := a$.
  \item $FD_3$ ($\at{ticker} \rightarrow S$): rows \texttt{stock}, \texttt{holds} and \texttt{transaction} agree on $\at{ticker} = a$, and \texttt{stock} has $S = a$. Set $b_{35} := a$ and $b_{45} := a$.
  \item $FD_5$ ($\at{user\_id}, \at{ticker} \rightarrow P$): rows \texttt{holds} and \texttt{transaction} agree on $\at{user\_id} = a$ and $\at{ticker} = a$, and \texttt{holds} has $P = a$. Set $b_{47} := a$.
\end{enumerate}

\textbf{Final tableau:}
\[
\begin{array}{l|ccccccc}
  & \at{transaction\_id} & \at{user\_id} & N & \at{ticker} & S & T & P \\
  \hline
  \texttt{user}        & b_{11} & a      & a      & b_{14} & b_{15} & b_{16} & b_{17} \\
  \texttt{stock}       & b_{21} & b_{22} & b_{23} & a      & a      & b_{26} & b_{27} \\
  \texttt{holds}       & b_{31} & a      & a      & a      & a      & b_{36} & a      \\
  \texttt{transaction} & a      & a      & a      & a      & a      & a      & a      \\
\end{array}
\]

The \texttt{transaction} row is now all distinguished symbols ($a$), so the decomposition is \textbf{lossless}. This agrees with the step-by-step proof.

\subsection*{Joining the Implemented Tables}

The proof covers instances of $U$, where every user, stock and position belongs to at least one trade. The implemented tables can store more than $U$ can, which is how they remove the insertion anomalies. For example, \texttt{user} stores Isla MacDonald, who has no trades.

Tuples that have no matching tuple in another relation are \textbf{dangling tuples}, and the natural join drops them. In the seed data:
\begin{itemize}[leftmargin=*]
  \item user 9 and the stocks CNR and ATD join with no trades;
  \item 16 of the 82 trades are in positions that have since been closed, and the seed data has no \texttt{holds} row for a closed position. The join of all four tables therefore returns 66 rows, one for each trade in an open position.
\end{itemize}
Dangling tuples are data that $U$ could never have held. They do not contradict losslessness: projecting any legal instance of $U$ and joining the projections back still returns exactly that instance.

\section{Counterexample: A Lossy Decomposition}

To show why the test matters, decompose \texttt{holds} so that the shared attributes are not a superkey of either side:
\[
  H_1 = \{\at{user\_id}, \at{h\_quantity}\}, \qquad H_2 = \{\at{ticker}, \at{h\_quantity}, \at{average\_price}\}
\]
Then $H_1 \cap H_2 = \{\at{h\_quantity}\}$, and $\cl{\at{h\_quantity}} = \{\at{h\_quantity}\}$, so neither $\{\at{h\_quantity}\} \rightarrow H_1$ nor $\{\at{h\_quantity}\} \rightarrow H_2$ holds. The test fails.

The seed data has two positions of 33 shares, held by different users in different stocks:
\begin{center}
\begin{tabular}{c c c c}
  \toprule
  \texttt{user\_id} & \texttt{ticker} & \texttt{h\_quantity} & \texttt{average\_price} \\
  \midrule
  3 & AMZN & 33 & 199.6300 \\
  8 & AAPL & 33 & 226.1091 \\
  \bottomrule
\end{tabular}
\end{center}

$\pi_{H_1}$ contains $(3, 33)$ and $(8, 33)$, and $\pi_{H_2}$ contains $(\text{AMZN}, 33, 199.6300)$ and $(\text{AAPL}, 33, 226.1091)$. Joining them on $\at{h\_quantity} = 33$ pairs every user with every stock:
\begin{center}
\begin{tabular}{c c c c l}
  \toprule
  \texttt{user\_id} & \texttt{ticker} & \texttt{h\_quantity} & \texttt{average\_price} & \\
  \midrule
  3 & AMZN & 33 & 199.6300 & original \\
  3 & AAPL & 33 & 226.1091 & \textbf{spurious} \\
  8 & AMZN & 33 & 199.6300 & \textbf{spurious} \\
  8 & AAPL & 33 & 226.1091 & original \\
  \bottomrule
\end{tabular}
\end{center}

Over all 22 rows of \texttt{holds}, $\pi_{H_1}(\texttt{holds}) \bowtie \pi_{H_2}(\texttt{holds})$ has 24 rows: the 22 originals plus these 2 spurious tuples. The join claims that user 3 holds AAPL and user 8 holds AMZN, which is false. Since the original relation cannot be recovered, the decomposition is \textbf{lossy}:
\[
  \pi_{H_1}(\texttt{holds}) \bowtie \pi_{H_2}(\texttt{holds}) \supsetneq \texttt{holds}
\]

\end{document}
